\subsection{4.2 扩展 RL 环境与脚手架}\label{scaling-rl-environments-and-harnesses}

扩展智能体强化学习需要多样化的环境：这些环境要能反映真实世界任务、支持训练规模下的可复现执行，并提供与任务目标一致的奖励。同时满足这些要求并不容易：真实工作流涉及异构工具与复杂状态，而自动评分器可能不完整、不一致，或者容易被钻空子。因此我们投入大规模的环境合成与整理，覆盖代码（§4.2.1）、通用专业工作流（§4.2.2）、视觉产物（§4.2.3）和网络安全（§4.2.4），并强调监督质量。本节介绍我们在这些领域如何构造任务与环境、智能体借以与环境交互的脚手架（§4.2.5），以及用于提升奖励可靠性、在训练前和训练中缓解奖励作弊的执行检查、基于 rollout 的审计和对抗性测试（§4.2.6）。

\subsection{4.2.1 代码智能体任务}\label{code-agent-tasks}

代码任务特别适合 RL，因为候选解可以通过自动化测试执行并评估，从而在规模化下提供反馈。然而，仅靠可执行的评估并不能保证监督可靠：测试可能遗漏必需的行为、拒掉合法实现，或者在不同次执行中给出不一致的结果。核心挑战因此是：在扩展任务构造规模的同时，保证奖励如实反映任务的正确性。为此，我们投入大量精力为代码智能体任务构建可扩展的合成与整理流水线，重点保证监督的准确性与稳健性。图 4 概括了任务合成路径以及用于评估监督准确性与稳健性的各类检查。

从多样来源进行可扩展的任务合成 我们通过五条互补的合成路径扩展任务构造，覆盖异构的编程语言、开发场景和任务时间跨度。第一，在基于 GitHub 的合成中，我们收集 pull request 及其关联的 issue，并用启发式和模型辅助的过滤器剔除重复、畸形或不可复现的候选。当某个 pull request 关联的 issue 对预期改动给出了足够完整的描述时，我们就把该 issue 用作任务规格。否则，由 LLM 依据参考补丁和周边仓库上下文重建一份 issue 式的规格，同时明确指示它省略可能泄露解法的实现细节。第二，为覆盖日常开发工作流，包括交互式的「氛围编程」（vibe coding）场景，我们收集组织内员工贡献的真实开发需求。这些需求覆盖功能实现、重构、调试、仓库维护等常见活动，智能体需要为其生成并实现用例，把所需行为落地为可执行的形式。第三，规格驱动的任务强调在详细、多约束的要求下生成代码，用于评估模型能否把复杂规格忠实转化为可执行的实现。第四，在源码驱动的合成中，我们使用 CodeMidas（Ye et al., 2026）从既有代码库已实现的功能出发派生任务。智能体识别候选功能，并将其可观测行为转化为任务规格与可执行环境；这条路径无需 issue、pull request 或其他开发产物，即可跨多样仓库扩展任务构造。对于长跨度软件工程任务，智能体会迭代细化任务需求、扩大所需改动的范围、引入跨组件依赖，并构造测试来检验由此产生的多步行为。我们还纳入并过滤来自公开来源（Badertdinov et al., 2026; PrimeIntellect, 2026; Tao et al., 2026; Yang et al., 2026b; Zan et al., 2025; Zhao et al., 2026）和授权数据供应商的样本，以进一步丰富任务分布。这些路径共同产出的任务覆盖广泛的编程语言、开发语境、实现复杂度和任务时间跨度。




\begin{figure}[htbp]
\centering
\includegraphics[width=\linewidth]{images_hi/fig04.png}
\caption{代码智能体任务扩展流水线总览。我们从多样来源构建任务，并经过严格的评估阶段，以确保监督准确且稳健。}
\label{fig:4}
\end{figure}

监督的准确性 我们让每个任务规格的行为范围与其单元测试的行为范围保持一致。目标是接受满足规格的实现、拒绝违反规格的实现，而不要求参考实现中那些无关紧要的细节。我们会审查任务是否存在覆盖不足和检查过于严苛的问题。尤其是，对于那些强制了规格中并不存在的要求的测试，我们会加以修订或删除。由于仅靠人工检查可能漏掉规格与测试之间的不一致，我们进一步通过基于 rollout 的审计来评估二者的对齐程度。每个任务由代码智能体尝试四次。审计智能体会收到全部四次 rollout，置于一个共享工作区中，其中包含问题陈述、测试、参考补丁，以及每次 rollout 提交的补丁、测试输出和完整对话日志。审计器首先阐明正确解需要满足什么，并评估测试是否覆盖了这些要求。随后，它借助补丁和对话日志，按规格逐一评估每个提交的解，并将这一评估与观测到的奖励作对比。被判定为不正确却通过了的解会被标记为潜在假阳性，说明验证不完整。被判定为正确却失败了的解会被标记为潜在假阴性，说明测试过于严苛或执行失败。这些分歧能指出需要进一步审查的任务，并为预期行为与所实现的验证之间可能存在的失配提供具体证据。

监督的稳健性 我们通过重复执行来检查奖励的稳定性。对于带有参考补丁的任务，我们要求 fail-to-pass（F2P）测试在应用参考补丁前必须失败、pass-to-pass（P2P）测试在应用参考补丁前必须通过；应用补丁后，两者都必须通过。我们要求这些结果在八次重跑中保持稳定，以此筛除不稳定测试和环境引起的奖励波动。我们还复用此前的 rollout 流水线收集智能体轨迹，以考察潜在的奖励作弊。提交的补丁和对话日志提供了智能体如何获得奖励的证据，有助于识别非预期的捷径，或对环境与评估流程的利用。这些检查补充了对解正确性的评估，并为下文所述的更广泛的奖励作弊缓解流程提供依据。

上述流程共同支撑在统一的监督质量标准下构建多样的代码任务语料。规格—测试对齐针对的是奖励信号的保真度，重复执行检查的是其稳定性，基于 rollout 的审计考察的是它与独立的解正确性评估是否一致。我们的目标是让这一组合可扩展，使 RL 训练既能受益于广泛的任务覆盖，又能引导智能体满足预期要求，而不是去利用评估中的漏洞。

\subsection{4.2.2 通用智能体任务}\label{general-agent-tasks}

我们把智能体能力扩展到跨领域的真实专业工作流，这些工作流要求智能体分析异构文档、使用专业软件，并产出符合领域标准的交付物。我们合成支持此类工作流的环境，再构造带有可验证结果的挑战性任务。图 5 概括了这条流水线，包括环境构建、任务合成和迭代式评分细则精炼。

环境设计 通用智能体环境通常由一个工作区和一组软件工具构成，通过 MCP、API、CLI 和 GUI 等接口访问。我们的设计在真实性与大规模 RL 的要求之间取得平衡：文件、软件行为和底层数据应当反映专业实践，同时执行必须保持在本地，且环境要易于重置。因此我们收集真实文件，直接纳入或作为合成时的参考，并使用自动化的多智能体工作流构建本地软件 mock，复现其支持的操作、状态变化、输出格式和错误响应。本身不依赖外部服务即可运行的工具则直接集成。所有状态都存储在本地，每次 rollout 都在一个可恢复到固定初始状态的隔离沙箱中运行。这避免了网络波动和服务限制，同时支持可复现执行。




\begin{figure}[htbp]
\centering
\includegraphics[width=\linewidth]{images_hi/fig05.png}
\caption{通用智能体环境与可验证任务合成流水线总览。我们首先基于真实文件与软件 mock 构建真实、可重置的环境。随后智能体探索这些环境，并合成带有可验证评分细则的任务。}
\label{fig:5}
\end{figure}

环境合成 为把这些组件组装成连贯的场景，规划智能体负责指定工作区结构、选择合适的工具，并规划文件和数据库内容。网络搜索让规划以真实世界信息为依据，而文件之间以及文件与数据库记录之间的关系则支撑后续的多跳推理与交叉验证。随后，多个智能体按照共享的规划以及各 mock 的数据规格，并行生成文件并填充数据库。审查智能体检查各个产物以及全局一致性，包括实体名称、数值核对、时间线和引用。任务构造之前，通过迭代修复解决不一致之处。工作区和工具配置随场景而异，因此每个环境只反映相关的工作实践，而不必具备所有可用的文件类型或软件接口。

任务合成与验证 我们调研常见的专业任务，将其组合并泛化为种子任务。智能体探索每个环境，并借助这些种子构造与其内容和工具相适配的任务。任务完成情况使用原子化、二元的评分细则条目来评估：基于代码的检查验证确定性属性，例如数据库取值和交付物格式；基于 LLM 的检查则评估更开放式的内容。对于基于 LLM 的条目，同一模型重复评判之间以及不同评判模型之间的一致性有助于识别歧义。

一致性本身并不能说明评分细则捕捉到了任务完成情况。因此我们收集能力水平不同的模型产生的 rollout，并让审查智能体考察它们的轨迹、执行结果和评分细则判定。审查者会修订那些拒掉合法解的过严标准，以及那些接受不完整解的过宽标准。为检验对奖励作弊的抵抗能力，我们加入针对无关文件或数据库被非预期改动的负向检查，并构造看似正确但并未真正解决任务的对抗性解。在 RL 期间，由自托管的 MiMo-V2.6-SFT 模型充当评分器，以支持稳定的评分。这一流程产出数千个环境和多样、可验证的任务，我们将其与其他任务合并用于 RL 训练。
